\scanpage{52}
% archive sheet number: 51
Algèbres étales et la différente \mspage{(4)}\hfill Tribu n\textsuperscript{o} 65

\S~2. \uline{Différente}.

\uline{n\textsuperscript{o} 1}. \uline{Idéaux de Filting d'un module} $M$.

Voir thèse de D. Lazard. Lien avec le contenu. Le premier Filting est compris entre
$\operatorname{Ann}(M)$ et une de ses puissances.

\uline{n\textsuperscript{o} 2}. \uline{Différente de Kähler}.

On prend pour $B$ une algèbre de p. f. sur $A$ (plus généralement il suffit peut
être de supposer que $\handwritten{\Omega}_{A}(B)$ est un $B$-module de p. f.). Par définition la différente
de Kähler $\text{\uline{$D$}}_{B/A}$ est le premier idéal de Filting de $\handwritten{\Omega}_{A}(B)$ :
\[
\text{si }\quad A^{p}\xrightarrow{u}A^{q}\longrightarrow\handwritten{\Omega}_{A}(B)\longrightarrow 0,\quad \text{c'est }\operatorname{Im}(\handwritten{\unsure{\Lambda|\bigwedge}}^{q}u).
\]
La variété des zéros de $\text{\uline{$D$}}_{B/A}$ est l'ensemble des points de ramification. Comportement
de $\text{\uline{$D$}}_{B/A}$ par changement d'anneau (c'est trivial).

\uline{n\textsuperscript{o} 3}. \uline{Transitivité}.

Dire ce qu'on sait. Si on a $A\to B\to C$ avec $C$ intersection complète locale
sur $B$, la formule de transitivité marche.

\uline{n\textsuperscript{o} 4}. \uline{Lien avec la trace}.

Définir la différente de Dedekind par la trace, et montrer qu'elle coïncide avec
celle de Kähler quand $A$ et $B$ sont des anneaux de valuation discrète (il s'agit
du gros th. 3 du n\textsuperscript{o} 7, \S~2, avec son dévissage, - éventuellement agrémenté des lemmes
du n\textsuperscript{o} 6 s'ils n'ont pas été faits dans un chapitre antérieur pour illustrer la notion
d'intersection complète. En tous cas on aura dit que si $A$ et $B$ sont réguliers,
$A\longrightarrow B$ est une intersection complète relative. La prop. 11 sera éclairée par le lemme
suivant :

``Si $k\longrightarrow R$ est finie et si $R/\operatorname{rad}(R)$ est étale, alors $R$ monogène ssi les
idéaux maximaux de $R$ sont homogènes''.

La prop. 12 est absorbée par le résultat ci-dessus sur les anneaux réguliers.

% Unlabelled faint outline in the left margin.
\figure{\begin{tikzpicture}[x=1mm,y=1mm]
\draw[gray,thin] (2.3,4.2) .. controls (1.9,4.9) and (1.1,4.5) .. (1.3,4.0)
 .. controls (1.5,3.7) and (1.9,3.8) .. (2.0,4.0)
 .. controls (1.5,3.7) and (.7,3.8) .. (.4,4.0)
 .. controls (.1,4.8) and (.8,5.8) .. (1.6,5.8)
 .. controls (2.2,5.8) and (2.3,5.1) .. (2.8,5.7)
 .. controls (3.4,6.6) and (3.5,6.5) .. (4.0,6.0)
 .. controls (4.4,5.7) and (4.7,6.4) .. (4.9,6.6)
 .. controls (5.1,6.8) and (5.1,6.2) .. (5.5,6.3)
 .. controls (5.8,6.8) and (6.1,5.8) .. (6.2,5.5)
 .. controls (6.4,5.1) and (6.2,4.7) .. (5.9,4.4)
 .. controls (6.7,4.6) and (7.2,4.4) .. (7.3,3.7)
 .. controls (7.5,3.9) and (7.8,3.3) .. (7.9,3.1)
 .. controls (8.4,2.4) and (8.2,1.0) .. (8.7,.4);
\end{tikzpicture}}

\uline{n\textsuperscript{o} 5}. \uline{Compléments}.

La ramification modérée du n\textsuperscript{o} 5 (noter que la formule de réduction des traces
utilisée en haut de la p. 25 est faite en Alg. VIII, \S~11).

Donner la description Kummerienne des extensions totalement et modérément ramifiées
(extraire la racine $n$-ième d'une uniformisante bien choisie ; ceci dans le cas complet);
exemple des surfaces de Riemann (ici le groupe de Galois a un générateur canonique
$\exp\left(\dfrac{2i\handwritten{\pi}}{n}\right)$ ).

